\chapter{Checkpoint \& Resume}
\label{ch:checkpoint}

\newthought{Scans get interrupted} --- power, quotas, preemption, your own Ctrl+C. \Jtwo\
checkpoints the control process every thirty seconds and can resume a scan with
\emph{no missing and no duplicated points}. This chapter explains what a checkpoint
contains, the safe-barrier rule that makes resume exact, and the operational details.

\section{What is saved, and what deliberately is not}
\label{sec:checkpoint-contents}

The Sampler is the source of truth. A checkpoint stores:

\begin{itemize}
  \item the \textbf{sampler's runtime state} --- RNG streams (as serialized seed
        sequences), the generated point set or file cursor, and per-method internals;
  \item the \textbf{submitted and completed \textsc{uuid} sets};
  \item a \textbf{run spec}: the normalized configuration with its hash, scan/task
        metadata, and the variable signature --- so a resume against a \emph{changed}
        card is detected rather than silently mixed.
\end{itemize}

\begin{jnote}
\textbf{In-flight tasks are never serialized.} Redis queues are volatile by design: on
resume, stale queue entries are drained and calculator pools rebuilt, and the sampler
simply re-proposes whatever was submitted but never acknowledged by the Archiver.
Combined with idempotent archiving, the arithmetic closes: no lost points, no doubles.
\end{jnote}

\section{The safe barrier}
\label{sec:safe-barrier}

A checkpoint is written only at a \emph{safe barrier}, defined by two conditions:

\begin{enumerate}
  \item the sampler is at rest (not mid-generation, not mid-batch);
  \item every submitted \textsc{uuid} has been acknowledged by the Archiver.
\end{enumerate}

A background heartbeat attempts a save every \textbf{30 seconds}; ticks that miss the
barrier are simply skipped, so the file on disk is always a consistent cut of the run.
Writes are atomic (temp file + rename) --- a crash mid-save cannot corrupt the previous
checkpoint.

\section{Where it lives}
\label{sec:checkpoint-location}

\begin{quote}
\file{<project-root>/checkpoints/<scan>/<sampler>/state.pkl}
\end{quote}

One file per scan and sampler, overwritten in place. Checkpoints are project-local like
everything else; deleting the directory is the supported way to force a fresh start.

\section{Resuming}
\label{sec:resuming}

Run the same card again:

\begin{terminal}
\begin{jcmd}
Jarvis2 run scan.yaml            # prompts when a checkpoint exists
Jarvis2 run scan.yaml --resume   # accepts it non-interactively
\end{jcmd}
\end{terminal}

On accept, \Jtwo\ validates the checkpoint (format, version, config hash), restores the
sampler state, drains any stale Redis queue, rebuilds calculator pools, and re-proposes
the unacknowledged points. Because every sampler mints deterministic \textsc{uuid}s and
per-Sample seeds are derived from a master seed sequence --- independent of Worker count
--- a resume may use a \emph{different} number of Workers, or a different machine, and
still produce the identical point set.

\begin{jwarn}
A checkpoint whose format is not the current distributed-runtime format --- including
files produced by other \Jarvis\ product lines --- is \textbf{refused} with an
actionable message, never partially loaded. Likewise, editing the card between runs
changes the config hash; resume tells you instead of continuing a subtly different
scan. Change of plan = new scan name.
\end{jwarn}

\section{Failure recovery during a run}
\label{sec:failure-recovery}

Resume handles the control process dying. The runtime also self-heals below that
level, with no operator action:

\begin{itemize}
  \item \textbf{A Worker dies} (OOM, node fault): its heartbeat lapses, the factory
        respawns it, and the Sample it held is re-queued.
  \item \textbf{A Sample fails} (calculator crash, timeout, missing observable): it is
        archived as a failed record with its log; the scan continues.
  \item \textbf{The Archiver falls behind}: results queue in Redis; Workers keep
        computing; buckets pack only after records land
        (Section~\ref{sec:sample-directory}).
\end{itemize}

\begin{jtip}
For long scans on shared machines, make \cli{Jarvis2 run scan.yaml --resume} the
\emph{only} way you launch: first invocation starts fresh (no checkpoint exists),
every later invocation continues. Your batch-system retry hook is then a one-liner.
\end{jtip}
